\chapter{Testing, Vibe Coding, and the Final Project}\label{chap:testing}
\courseweek{Weeks 15--18.} Vibe coding is a fast way to express an intent to an AI coding tool. It does not transfer accountability. The course uses it as an implementation technique within a testable engineering workflow.

\begin{learningbox}
\begin{itemize}
  \item Turn a natural-language change request into observable acceptance criteria.
  \item Review an AI-generated patch for data flow, assumptions, and failure behaviour.
  \item Design a small test suite before or alongside an agent implementation.
  \item Demonstrate a final workflow with evidence, limits, and a rollback path.
\end{itemize}
\end{learningbox}

\section{Vibe coding with boundaries}
A responsible request names the intended user, input/output contract, constraints, and tests. After a tool proposes code, inspect the diff, run the relevant notebook cells or tests, and explain what changed. Avoid asking a model to make broad changes to an unexamined codebase. Prefer small, reversible changes with a visible before/after result.

\section{Evaluation before demonstration}
The final project needs a compact evaluation set that includes ordinary cases, edge cases, and known failure cases. For each case, record expected behaviour, actual behaviour, evidence, and follow-up action. A polished demo without a reproducible trace does not clear the project requirement.

\section{Final assessment}
The final project demonstration explains the problem boundary, agent or workflow design, evidence from tests, remaining limitations, and a next safe step. The written examination checks conceptual transfer. The oral examination can ask students to reason about a tool contract, an AI-generated patch, an evaluation trace, a rollback, or an unexpected failure.

\begin{labbox}
\textbf{Weeks 15--17:} use the project studio to build in small steps. Preserve the baseline, record every meaningful change, and rehearse the explanation of one failure case before the final demonstration.
\end{labbox}

\section{Risk-aware evaluation}
For the final project, evaluate more than apparent task completion. Include an ordinary case, an adverse or ambiguous case, a tool failure, and a case that must escalate to a human. Record the risk, expected behaviour, observed behaviour, evidence, and mitigation. This course-level practice is consistent with the lifecycle emphasis of the NIST Generative AI Profile, while remaining small enough for a student project \citep{nistgenai2024}.
